\documentclass[10pt,a4paper]{amsart}
\usepackage[margin=19mm]{geometry}
\usepackage{amsmath,amssymb,mathtools}
\usepackage[scr=rsfs,cal=euler]{mathalfa}
\usepackage{xfrac}
\usepackage[unicode,hidelinks,bookmarks=false]{hyperref}
\newcommand{\ie}{i.e.\ }
\newcommand{\cf}{cf.\ }
\newcommand{\resp}{respectively\ }
\newcommand{\Subsection}{\subsection}
\providecommand{\nobreakdash}{\nobreak\hbox{-}}
\setlength{\emergencystretch}{2em}
\multlinegap0pt
\usepackage{tikz}
\newtheorem{theorem}{Theorem}
\newcommand{\B}[1]{{\mathbf #1}}
\newcommand{\F}[1]{{\mathfrak #1}}
\title{A refinement of bounded cohomology}
\author{Światosław Gal, Jarek Kędra}
\date{Source: arXiv:2208.03168v3 (CC BY 4.0)}
\begin{document}
\maketitle

\noindent\emph{Source and licence.} A refinement of bounded cohomology. Version arXiv:2208.03168v3 (CC BY 4.0), distributed by arXiv under the Creative Commons Attribution 4.0 International licence, http://creativecommons.org/licenses/by/4.0/, as recorded in the arXiv licence metadata for this exact version and shown on the abstract page https://arxiv.org/abs/2208.03168v3. Full licence notice: \texttt{licenses/arxiv-2208.03168-CC-BY-4.0.txt}.\par\smallskip

\noindent\textit{Editorial scope.} Counted unit: the article's theorem T:comparison (source0.tex lines 300--306). The article's complete original proof of it is delivered with it (source0.tex lines 308--345, one \texttt{proof} environment of 38 lines). No mathematics is written, repaired or re-derived: the statement and the proof are the article's own lines, in the article's own notation, and no retained line is substituted or re-flowed.  No further source-local context is needed: every cross-reference of every retained line resolves inside the two delivered spans.  Nothing was omitted from any retained span, so the package carries no omission record.  No retained line contains a citation, so the wrapper carries no bibliography and no bibliography entry.  No retained line contains a figure environment, an image-inclusion command or drawing code, so no image is delivered and the package declares no asset dependency.  Editorial formatting only, in addition: an AMS \texttt{amsart} preamble that restates the article's own declarations and macros used by the retained lines, namely the environments \texttt{theorem} on separate counters, together with the standard packages those lines need.  The retained environments print in this document as Theorem 1 in order of appearance, whereas the article prints the same items with the numbering of its own full document.  The complete native source of the article as distributed in the arXiv e-print for this exact version is bundled unchanged as \texttt{sources/129137/source0.tex}.
\begin{theorem}\label{T:comparison}
If\/ $\Gamma$ is an amenable group then the comparison homomorphisms
$$
H^n_{(p)}(\Gamma;\B R)\to H^n_{(p-1)}(\Gamma;\B R)
$$
vanish for all $n>0$ and $p=1,2,\dots, n$.
\end{theorem}

\begin{proof}
The argument follows the proof of the fact that bounded cohomology
of an amenable group vanishes in positive degree.
Let $\F m$ be a right-invariant mean on $\Gamma$.  Define a map
$M \colon C^{n+1}_{(p)}(\Gamma;\B R)\to C^n_{(p-1)}(\Gamma;\B R)$ by
$$
M\omega(g_1,\dots,g_n)=\int\omega(h,g_1,\dots,g_n)\,\F m(h).
$$
To see that $M\omega$ is $(p-1)$-bounded fix $g_p,\ldots,g_n$ and
observe that the function
$h\mapsto \omega(h,g_1,\ldots,g_{p-1},g_p,\ldots,g_n)$
is bounded uniformly for any choice of $g_1,\ldots,g_{p-1}$. 
This shows $(p-1)$-boundedness of $M\omega$.

Then
\begin{align*}
M\delta\omega(g_1,\dots,g_n)
&=\int\delta\omega(h,g_1,\dots,g_n)\,\F m(h)\\
&=\int\omega(g_1,\dots,g_n)\F m(h)\\
&\hphantom{=}-\int\omega(hg_1,g_2,\dots,g_n)\,\F m(h)+\int\omega(h,g_1g_2,\dots,g_n)\,\F m(h)-\dots\\
&\hphantom{=}\dots\pm \int\omega(h,g_1,\ldots,g_{n-1})\,\F m(h)\\
&\overset{!}=\omega(g_1,\dots,g_n)\\
&\hphantom{=}-\int\omega(h,g_2,\dots,g_n)\,\F m(h)+\int\omega(h,g_1g_2,\dots,g_n)\,\F m(h)-\dots\\
&\hphantom{=}\dots\pm \int\omega(h,g_1,\ldots,g_{n-1})\,\F m(h)\\
&=\omega(g_1,\dots,g_n)-M\omega(g_2,\dots,g_n)+M\omega(g_1g_2,\dots,g_n)-\dots\\
&\hphantom{=}\dots\pm M\omega(g_1,\ldots,g_{n-1})\\
&=\omega(g_1,\dots,g_n)-\delta M\omega(g_1,\dots,g_n),
\end{align*}
where, in the marked equality, we used the invariance of the mean 
in the second term.
We obtain that
$$
M\delta\omega+\delta M\omega=\omega=i(\omega),
$$
where $i\colon C^n_{(p)}(\Gamma;\B R)\to C^n_{(p-1)}(\Gamma;\B R)$ is the inclusion
of $p$-bounded cochains into $(p-1)$-bounded ones.
Consequently, the map $M$ is a homotopy for the comparison map.
\end{proof}

\end{document}
